\documentclass[11pt]{article}
\usepackage[utf8]{inputenc}
\usepackage[T1]{fontenc}
\usepackage{multicol}
\usepackage{geometry}
\usepackage{enumitem}
\usepackage{xcolor}
\usepackage{titlesec}
\usepackage{ragged2e}
\usepackage{amsmath}

\geometry{a4paper, left=1cm, right=1cm, top=0.8cm, bottom=1.2cm}

\setlength{\parindent}{0pt}
\setlength{\columnsep}{1cm}
\setlength{\columnseprule}{0.4pt}

\titleformat{\section}{\normalfont\Large\bfseries\color{blue}}{}{0em}{}

\title{\textbf{\textcolor{red}{Juros Simples - Correção}}}
\author{Professor: Jefferson}
\date{}

\begin{document}

\maketitle
\vspace{-0.5cm}

\begin{multicols}{2}

\section*{Questões e Respostas}

\begin{enumerate}[leftmargin=*, itemsep=2pt]

\item Calcule os juros simples: $C = R\$ 1000$, $i = 2\%$ a.m., $t = 6$ meses \\
\textcolor{blue}{\textbf{Resposta: R\$ 120,00}} \\
\textcolor{blue}{\textbf{Resolução:} $J = C \cdot i \cdot t = 1000 \cdot 0,02 \cdot 6 = 120$}

\item Calcule os juros simples: $C = R\$ 500$, $i = 1,5\%$ a.m., $t = 1$ ano \\
\textcolor{blue}{\textbf{Resposta: R\$ 90,00}} \\
\textcolor{blue}{\textbf{Resolução:} Converter: $t = 12$ meses. $J = 500 \cdot 0,015 \cdot 12 = 90$}

\item Calcule os juros simples: $C = R\$ 2000$, $i = 10\%$ a.a., $t = 2$ anos \\
\textcolor{blue}{\textbf{Resposta: R\$ 400,00}} \\
\textcolor{blue}{\textbf{Resolução:} $J = 2000 \cdot 0,10 \cdot 2 = 400$}

\item Calcule o montante: $C = R\$ 800$, $i = 3\%$ a.m., $t = 5$ meses \\
\textcolor{blue}{\textbf{Resposta: R\$ 920,00}} \\
\textcolor{blue}{\textbf{Resolução:} $M = C(1 + i \cdot t) = 800 \cdot (1 + 0,03 \cdot 5) = 800 \cdot 1,15 = 920$}

\item Calcule o montante: $C = R\$ 1500$, $i = 2\%$ a.m., $t = 8$ meses \\
\textcolor{blue}{\textbf{Resposta: R\$ 1.740,00}} \\
\textcolor{blue}{\textbf{Resolução:} $M = 1500 \cdot (1 + 0,02 \cdot 8) = 1500 \cdot 1,16 = 1740$}

\item Calcule o montante: $C = R\$ 3000$, $i = 15\%$ a.a., $t = 3$ anos \\
\textcolor{blue}{\textbf{Resposta: R\$ 4.350,00}} \\
\textcolor{blue}{\textbf{Resolução:} $M = 3000 \cdot (1 + 0,15 \cdot 3) = 3000 \cdot 1,45 = 4350$}

\item Complete a tabela: $C = R\$ 1000$, $i = 2\%$ a.m., $t = 6$ meses, $J = ?$ \\
\textcolor{blue}{\textbf{Resposta: R\$ 120,00}} \\
\textcolor{blue}{\textbf{Resolução:} $J = 1000 \cdot 0,02 \cdot 6 = 120$}

\item Complete a tabela: $C = R\$ 500$, $i = 12\%$ a.a., $t = 2$ anos, $J = ?$ \\
\textcolor{blue}{\textbf{Resposta: R\$ 120,00}} \\
\textcolor{blue}{\textbf{Resolução:} $J = 500 \cdot 0,12 \cdot 2 = 120$}

\item Complete a tabela: $C = R\$ 2000$, $i = 1,5\%$ a.m., $t = 1$ ano, $J = ?$ \\
\textcolor{blue}{\textbf{Resposta: R\$ 360,00}} \\
\textcolor{blue}{\textbf{Resolução:} Converter: $t = 12$ meses. $J = 2000 \cdot 0,015 \cdot 12 = 360$}

\item João aplicou R\$ 800 a juros simples de 3\% a.m. por 5 meses. Qual o montante? \\
\textcolor{blue}{\textbf{Resposta: R\$ 920,00}} \\
\textcolor{blue}{\textbf{Resolução:} $M = 800 \cdot (1 + 0,03 \cdot 5) = 800 \cdot 1,15 = 920$}

\item Um capital de R\$ 1200 rendeu R\$ 288 em 1 ano. Qual a taxa anual? \\
\textcolor{blue}{\textbf{Resposta: 24\% a.a.}} \\
\textcolor{blue}{\textbf{Resolução:} $i = \frac{J}{C \cdot t} = \frac{288}{1200 \cdot 1} = 0,24 = 24\%$}

\item Qual o tempo necessário para R\$ 2000 render R\$ 600 a 5\% a.m.? \\
\textcolor{blue}{\textbf{Resposta: 6 meses}} \\
\textcolor{blue}{\textbf{Resolução:} $t = \frac{J}{C \cdot i} = \frac{600}{2000 \cdot 0,05} = \frac{600}{100} = 6$}

\item Converta 12\% a.a. para taxa mensal \\
\textcolor{blue}{\textbf{Resposta: 1\% a.m.}} \\
\textcolor{blue}{\textbf{Resolução:} $i_{\text{mensal}} = \frac{12\%}{12} = 1\%$}

\item Converta 24\% a.a. para taxa mensal \\
\textcolor{blue}{\textbf{Resposta: 2\% a.m.}} \\
\textcolor{blue}{\textbf{Resolução:} $i_{\text{mensal}} = \frac{24\%}{12} = 2\%$}

\item Converta 30\% a.a. para taxa mensal \\
\textcolor{blue}{\textbf{Resposta: 2,5\% a.m.}} \\
\textcolor{blue}{\textbf{Resolução:} $i_{\text{mensal}} = \frac{30\%}{12} = 2,5\%$}

\item Um capital de R\$ 2000 rendeu R\$ 600 em 1 ano. Qual a taxa anual? \\
\textcolor{blue}{\textbf{Resposta: 30\% a.a.}} \\
\textcolor{blue}{\textbf{Resolução:} $i = \frac{600}{2000 \cdot 1} = 0,30 = 30\%$}

\item Um capital de R\$ 2000 rendeu R\$ 450 em 9 meses. Qual a taxa mensal? \\
\textcolor{blue}{\textbf{Resposta: 2,5\% a.m.}} \\
\textcolor{blue}{\textbf{Resolução:} $i = \frac{450}{2000 \cdot 9} = \frac{450}{18000} = 0,025 = 2,5\%$}

\item Um capital de R\$ 2000 aplicado a 8\% a.a. por 2 anos. Qual o montante? \\
\textcolor{blue}{\textbf{Resposta: R\$ 2.320,00}} \\
\textcolor{blue}{\textbf{Resolução:} $M = 2000 \cdot (1 + 0,08 \cdot 2) = 2000 \cdot 1,16 = 2320$}

\item Em quanto tempo um capital duplica a 10\% a.a.? \\
\textcolor{blue}{\textbf{Resposta: 10 anos}} \\
\textcolor{blue}{\textbf{Resolução:} $2C = C(1 + 0,10 \cdot t) \Rightarrow 2 = 1 + 0,10t \Rightarrow t = 10$}

\item Em quanto tempo um capital triplica a 5\% a.m.? \\
\textcolor{blue}{\textbf{Resposta: 40 meses}} \\
\textcolor{blue}{\textbf{Resolução:} $3C = C(1 + 0,05 \cdot t) \Rightarrow 3 = 1 + 0,05t \Rightarrow t = 40$}

\item Em quanto tempo R\$ 1500 rende R\$ 600 a 2\% a.m.? \\
\textcolor{blue}{\textbf{Resposta: 20 meses}} \\
\textcolor{blue}{\textbf{Resolução:} $t = \frac{600}{1500 \cdot 0,02} = \frac{600}{30} = 20$}

\item Empréstimo de R\$ 5000 a 2\% a.m. por 10 meses. Qual o total a pagar? \\
\textcolor{blue}{\textbf{Resposta: R\$ 6.000,00}} \\
\textcolor{blue}{\textbf{Resolução:} $M = 5000 \cdot (1 + 0,02 \cdot 10) = 5000 \cdot 1,20 = 6000$}

\item Uma aplicação de R\$ 3000 rendeu R\$ 450 em 6 meses. Qual a taxa mensal? \\
\textcolor{blue}{\textbf{Resposta: 2,5\% a.m.}} \\
\textcolor{blue}{\textbf{Resolução:} $i = \frac{450}{3000 \cdot 6} = \frac{450}{18000} = 0,025 = 2,5\%$}

\item Quantos meses para R\$ 2000 virar R\$ 2600 a 1,5\% a.m.? \\
\textcolor{blue}{\textbf{Resposta: 20 meses}} \\
\textcolor{blue}{\textbf{Resolução:} $J = 2600 - 2000 = 600$. $t = \frac{600}{2000 \cdot 0,015} = 20$}

\item Do gráfico: Qual o capital inicial? \\
\textcolor{blue}{\textbf{Resposta: 2 (unidade do gráfico)}} \\
\textcolor{blue}{\textbf{Resolução:} Capital é o valor quando $t = 0$}

\item Do gráfico: Qual a taxa mensal? \\
\textcolor{blue}{\textbf{Resposta: 10\% a.m.}} \\
\textcolor{blue}{\textbf{Resolução:} Em 5 meses: $J = 3,0 - 2,0 = 1,0$. $i = \frac{1,0}{2 \cdot 5} = 0,10 = 10\%$}

\item Do gráfico: Qual o montante após 8 meses? \\
\textcolor{blue}{\textbf{Resposta: 3,6 (unidade do gráfico)}} \\
\textcolor{blue}{\textbf{Resolução:} $M = 2 \cdot (1 + 0,10 \cdot 8) = 2 \cdot 1,80 = 3,6$}

\item Um capital dobra em 10 anos a juros simples. Qual a taxa? \\
\textcolor{blue}{\textbf{Resposta: 10\% a.a.}} \\
\textcolor{blue}{\textbf{Resolução:} $2C = C(1 + i \cdot 10) \Rightarrow 2 = 1 + 10i \Rightarrow i = 0,10 = 10\%$}

\item Em quanto tempo um capital quadruplica a 5\% a.m.? \\
\textcolor{blue}{\textbf{Resposta: 60 meses}} \\
\textcolor{blue}{\textbf{Resolução:} $4C = C(1 + 0,05 \cdot t) \Rightarrow 4 = 1 + 0,05t \Rightarrow t = 60$}

\item Qual capital rende R\$ 1000 em 8 meses a 2,5\% a.m.? \\
\textcolor{blue}{\textbf{Resposta: R\$ 5.000,00}} \\
\textcolor{blue}{\textbf{Resolução:} $C = \frac{J}{i \cdot t} = \frac{1000}{0,025 \cdot 8} = \frac{1000}{0,20} = 5000$}

\item Banco A: R\$ 5000, 2,2\% a.m., 12 meses. Qual o juro? \\
\textcolor{blue}{\textbf{Resposta: R\$ 1.320,00}} \\
\textcolor{blue}{\textbf{Resolução:} $J = 5000 \cdot 0,022 \cdot 12 = 1320$}

\item Banco B: R\$ 5000, 2,5\% a.m., 10 meses. Qual o juro? \\
\textcolor{blue}{\textbf{Resposta: R\$ 1.250,00}} \\
\textcolor{blue}{\textbf{Resolução:} $J = 5000 \cdot 0,025 \cdot 10 = 1250$}

\item Qual banco cobra menos juros? \\
\textcolor{blue}{\textbf{Resposta: Banco B}} \\
\textcolor{blue}{\textbf{Resolução:} Banco B: R\$ 1.250,00 < Banco A: R\$ 1.320,00}

\item Um capital duplica em 8 anos. Em quanto tempo triplicaria? \\
\textcolor{blue}{\textbf{Resposta: 16 anos}} \\
\textcolor{blue}{\textbf{Resolução:} Se duplica em 8: $i = \frac{1}{8}$. Para triplicar: $3 = 1 + \frac{1}{8}t \Rightarrow t = 16$}

\item Se R\$ 4000 rende R\$ 1000 em 5 meses, qual a taxa anual? \\
\textcolor{blue}{\textbf{Resposta: 60\% a.a.}} \\
\textcolor{blue}{\textbf{Resolução:} Taxa mensal: $i_m = \frac{1000}{4000 \cdot 5} = 0,05 = 5\%$. Anual: $5\% \times 12 = 60\%$}

\item Depósitos mensais de R\$ 100 a 1\% a.m. por 12 meses. Montante total? \\
\textcolor{blue}{\textbf{Resposta: R\$ 1.278,00}} \\
\textcolor{blue}{\textbf{Resolução:} $M = 1200 + 100 \cdot 0,01 \cdot (12+11+\ldots+1) = 1200 + 78 = 1278$}

\item Compare: R\$ 1000 a 2\% a.m. por 2 anos vs R\$ 1000 a 1,8\% a.m. por 2 anos \\
\textcolor{blue}{\textbf{Resposta: 2\% a.m. rende R\$ 1.480,00; 1,8\% a.m. rende R\$ 1.432,00}} \\
\textcolor{blue}{\textbf{Resolução:} $M_1 = 1000(1 + 0,02 \cdot 24) = 1480$; $M_2 = 1000(1 + 0,018 \cdot 24) = 1432$}

\end{enumerate}

\end{multicols}

\end{document}
